\documentclass[11pt]{report}
\usepackage[margin=1in]{geometry}
\usepackage{amsmath,amssymb,mathtools,bm}
\usepackage{booktabs,array}
\usepackage{xcolor}
\usepackage{listings}
\usepackage[hidelinks]{hyperref}
\usepackage[capitalise,noabbrev]{cleveref}

\lstset{basicstyle=\ttfamily\small, columns=fullflexible, keepspaces=true,
        frame=single, framerule=0.3pt, rulecolor=\color{black!30},
        backgroundcolor=\color{black!3}, showstringspaces=false, language=Python,
        aboveskip=8pt, belowskip=8pt}
\numberwithin{equation}{chapter}
\allowdisplaybreaks

% unit vectors
\newcommand{\uh}{\hat{\mathbf u}}
\newcommand{\eh}{\hat{\mathbf e}}
\newcommand{\nh}{\hat{\mathbf n}}
\newcommand{\tth}{\hat{\mathbf t}}
\newcommand{\gh}{\hat{\mathbf g}}
\newcommand{\hh}{\hat{\mathbf h}}
\newcommand{\zh}{\hat{\mathbf z}_E}
\newcommand{\vb}[1]{\bm{#1}}
\newcommand{\pd}[2]{\frac{\partial #1}{\partial #2}}
\newcommand{\aL}{a_L}
\newcommand{\aD}{a_D}
\newcommand{\sdot}{\dot s}
\newcommand{\code}[1]{\texttt{#1}}

\title{Guidance and Control Models for the HL-20 Glide Vehicle\\[6pt]
\large Point-mass dynamics, linearization, error-state models for LQR,\\
and inner-loop attitude control (classical and NDI)}
\author{hgv3\_gc project notes}
\date{\today}

\begin{document}
\maketitle
\tableofcontents

%======================================================================
\chapter*{How to read this document}
\addcontentsline{toc}{chapter}{How to read this document}

This document derives, in long form, the models used by the trajectory planner and the
models needed to close the loop around the plan. The planner code is in
\code{waypoint\_traj\_solver.py} and \code{hl20\_model.py}.

The structure follows the pattern of Rajamani's
\emph{Vehicle Dynamics and Control}: one chapter derives the vehicle model from first
principles, and a later chapter goes back to the same nonlinear equations, substitutes
error states, and works them into a state-space model for LQR.

\begin{description}
\item[\Cref{ch:notation}] Frames, states, constants and sign conventions.
\item[\Cref{ch:eom}] The nonlinear point-mass equations of motion over a rotating,
  oblate Earth, derived from Newton's second law. The result is exactly the model in
  \code{hl20\_model.py}. The chapter ends with how it differs from the \code{hgv3}
  simulator.
\item[\Cref{ch:lin}] Jacobian linearization of those equations about a planned
  trajectory, every partial derivative written out.
\item[\Cref{ch:err}] The error-state (path-following) model: reference indexed by
  ground-track arc length, error substitution, and linearization into
  $\dot{\vb e} = A_e \vb e + B_e\,\delta\vb u$.
\item[\Cref{ch:lqr}] Using that model for LQR gain design, and computing the matrices
  numerically with CasADi.
\item[\Cref{ch:inner}] The inner (attitude) loop for a 6-DOF simulation: rigid-body
  rotational dynamics, wind-axis kinematics, classical gain-scheduled design, NDI and
  INDI, all in long form.
\item[\Cref{app:verify}] What was checked and how.
\end{description}

Equations are written component by component rather than in compact vector form, and
intermediate steps are kept. Every equation in Chapters 2--4 and the kinematics in
Chapter 6 were checked symbolically or numerically (\cref{app:verify}).

%======================================================================
\chapter{Notation, frames and conventions}\label{ch:notation}

\section{Reference frames}

\begin{description}
\item[Inertial (ECI), $I$.] Non-rotating, origin at Earth's center. The \code{hgv3}
  simulator integrates its equations of motion here.
\item[Earth-fixed (ECEF), $E$.] Rotates with Earth at constant rate $\Omega$ about its
  $z$ axis $\zh$. Positions and velocities in the planner are Earth-relative.
\item[Local geocentric frame, $(\uh,\eh,\nh)$.] At the vehicle's position: $\uh$ points
  radially up, $\eh$ east, $\nh$ north (perpendicular to $\uh$, toward the pole). It is
  right-handed in the order $(\uh,\eh,\nh)$:
  \[
    \uh\times\eh=\nh,\qquad \eh\times\nh=\uh,\qquad \nh\times\uh=\eh .
  \]
  It is \emph{geocentric}: $\uh$ is along the radius vector, not the ellipsoid normal.
\item[Velocity frame, $(\tth,\gh,\hh)$.] $\tth$ along the Earth-relative velocity,
  $\gh$ perpendicular to it in the vertical plane on the ``up'' side, and $\hh$
  horizontal and to the right of the velocity. Defined in \cref{sec:veldyn}.
\item[Body and wind axes.] Used only in \cref{ch:inner}.
\end{description}

\section{States, inputs and parameters}

\begin{center}
\begin{tabular}{@{}llll@{}}
\toprule
Symbol & Meaning & Units & Code name \\
\midrule
$r$ & geocentric radius & m & \code{r} \\
$\theta$ & longitude & rad & \code{lon} \\
$\phi$ & geocentric latitude & rad & \code{lat} \\
$V$ & Earth-relative speed & m/s & \code{V} \\
$\gamma$ & flight-path angle, positive climbing & rad & \code{gam} \\
$\psi$ & heading, clockwise from north & rad & \code{psi} \\
$\alpha$ & angle of attack & rad & \code{alpha} \\
$\sigma$ & bank angle, positive right wing down & rad & \code{sigma} \\
$u_\alpha,\ u_\sigma$ & $\dot\alpha,\ \dot\sigma$ (planner controls only) & rad/s & \code{U} \\
$h$ & geodetic altitude & m & \code{h} \\
$\aL,\ \aD$ & lift and drag \emph{accelerations} (force per unit mass) & m/s$^2$ & \code{L}, \code{D} \\
$\bar q$ & dynamic pressure $\tfrac12\rho V^2$ & Pa & \code{qbar} \\
$M$ & Mach number $V/a_s(h)$ & -- & \code{mach} \\
\bottomrule
\end{tabular}
\end{center}

The planner uses the 8-element state $[r,\theta,\phi,V,\gamma,\psi,\alpha,\sigma]$ with
$[u_\alpha,u_\sigma]$ as controls, so that $\alpha$ and $\sigma$ are smooth and
rate-limited. For the guidance problem in \cref{ch:lin,ch:err,ch:lqr}, $\alpha$ and
$\sigma$ are treated directly as the inputs, because the 3-DOF simulator flies them
exactly and a 6-DOF autopilot tracks them.

\section{Constants}

\begin{center}
\begin{tabular}{@{}llll@{}}
\toprule
Symbol & Value & Meaning & Source (\code{hgv3.}) \\
\midrule
$\mu$ & $3.986004418\times10^{14}$ m$^3$/s$^2$ & gravitational parameter & \code{MU\_EARTH} \\
$R_E$ & 6\,378\,137 m & equatorial radius & \code{R\_EARTH} \\
$J_2$ & $1.08263\times10^{-3}$ & oblateness coefficient & \code{J2} \\
$\Omega$ & $7.2921159\times10^{-5}$ rad/s & Earth rotation rate & \code{OMEGA\_EARTH} \\
$g_0$ & 9.80665 m/s$^2$ & standard gravity & \code{G0} \\
$a_{W}$ & 6\,378\,137 m & WGS84 semi-major axis & \code{WGS84\_A} \\
$f_{W}$ & $1/298.257223563$ & WGS84 flattening & \code{WGS84\_F} \\
$m$ & 8664 kg & HL-20 mass & airframe \\
$S$ & 26.61 m$^2$ & reference area & airframe \\
$\bar c$ & 8.61 m & reference length (pitch) & airframe \\
$b$ & 4.23 m & reference span (roll, yaw) & airframe \\
$I_{xx}, I_{yy}, I_{zz}$ & 10\,185, 45\,548, 48\,327 kg\,m$^2$ & principal inertias & airframe \\
$I_{xz}$ & 0 (assumed in the data) & product of inertia & airframe \\
\bottomrule
\end{tabular}
\end{center}

\section{Sign conventions}

\begin{itemize}
\item $\gamma>0$ is climbing; entry trajectories have $\gamma<0$.
\item $\psi$ is measured clockwise from north: $\psi=\pi/2$ is due east.
\item $\sigma>0$ (right wing down) tilts the lift vector to the right of the velocity and
  turns the vehicle right, so $\dot\psi$ increases. This matches \code{hgv3}.
\item Latitude in the planner is \emph{geocentric}. Waypoints and the simulator use
  \emph{geodetic} latitude; conversions go through ECEF (\code{hgv3.lla\_to\_ecef}).
\end{itemize}

%======================================================================
\chapter{Nonlinear point-mass equations of motion}\label{ch:eom}

The vehicle is modeled as a point mass acted on by aerodynamic force and gravity. There
is no thrust. Angle of attack $\alpha$ and bank $\sigma$ orient the aerodynamic force; their
own dynamics are the subject of \cref{ch:inner}.

\section{Newton's second law in the inertial frame}

With $\vb r$ the position vector from Earth's center, Newton's law in ECI is
\begin{equation}
  \left.\frac{d^2\vb r}{dt^2}\right|_I = \vb a_{\text{aero}} + \vb g(\vb r),
  \label{eq:newton}
\end{equation}
where $\vb a_{\text{aero}}$ is the aerodynamic force per unit mass and $\vb g$ is the
gravitational acceleration. This is the form the \code{hgv3} simulator integrates, with
RK4 on ECI position and velocity.

\section{Moving to the rotating Earth frame}

For any vector $\vb b$, the derivatives seen in the inertial and Earth frames are related
by the transport theorem,
\begin{equation}
  \left.\frac{d\vb b}{dt}\right|_I = \left.\frac{d\vb b}{dt}\right|_E + \vb\Omega\times\vb b,
  \qquad \vb\Omega = \Omega\,\zh .
\end{equation}
Applied to $\vb r$, it gives the inertial velocity in terms of the Earth-relative
velocity $\vb v \equiv \left.d\vb r/dt\right|_E$:
\begin{equation}
  \left.\frac{d\vb r}{dt}\right|_I = \vb v + \vb\Omega\times\vb r .
\end{equation}
Applying it a second time, with $\Omega$ constant:
\begin{align}
  \left.\frac{d^2\vb r}{dt^2}\right|_I
  &= \left.\frac{d}{dt}\big(\vb v + \vb\Omega\times\vb r\big)\right|_E
     + \vb\Omega\times\big(\vb v + \vb\Omega\times\vb r\big) \notag\\
  &= \left.\frac{d\vb v}{dt}\right|_E + \vb\Omega\times\vb v
     + \vb\Omega\times\vb v + \vb\Omega\times(\vb\Omega\times\vb r) \notag\\
  &= \left.\frac{d\vb v}{dt}\right|_E + 2\,\vb\Omega\times\vb v
     + \vb\Omega\times(\vb\Omega\times\vb r).
\end{align}
Substituting into \cref{eq:newton} and solving for the Earth-relative acceleration:
\begin{equation}
  \boxed{\;
  \left.\frac{d\vb v}{dt}\right|_E
  = \vb a_{\text{aero}} + \vb g
    \;\underbrace{-\,2\,\vb\Omega\times\vb v}_{\text{Coriolis}}
    \;\underbrace{-\,\vb\Omega\times(\vb\Omega\times\vb r)}_{\text{centripetal}} \;}
  \label{eq:rotating}
\end{equation}

\section{Spherical coordinates and the local basis}

The position is described by radius $r$, longitude $\theta$ and geocentric latitude $\phi$:
$\vb r = r\,\uh$. In ECEF components the local unit vectors are
\begin{align}
  \uh &= \begin{bmatrix} \cos\phi\cos\theta \\ \cos\phi\sin\theta \\ \sin\phi \end{bmatrix},&
  \eh &= \begin{bmatrix} -\sin\theta \\ \cos\theta \\ 0 \end{bmatrix},&
  \nh &= \begin{bmatrix} -\sin\phi\cos\theta \\ -\sin\phi\sin\theta \\ \cos\phi \end{bmatrix}.
\end{align}
Earth's spin axis, written in the local basis, is
\begin{equation}
  \zh = \sin\phi\,\uh + \cos\phi\,\nh,
  \qquad\text{so}\qquad
  \vb\Omega = \Omega\sin\phi\,\uh + \Omega\cos\phi\,\nh .
  \label{eq:omega_local}
\end{equation}

As the vehicle moves, $\theta$ and $\phi$ change and the local basis turns. Differentiating
the component expressions above with respect to time (in the Earth frame) gives
\begin{align}
  \frac{d\uh}{dt} &= \dot\phi\,\nh + \dot\theta\cos\phi\,\eh, \label{eq:duh}\\
  \frac{d\eh}{dt} &= -\dot\theta\cos\phi\,\uh + \dot\theta\sin\phi\,\nh, \\
  \frac{d\nh}{dt} &= -\dot\phi\,\uh - \dot\theta\sin\phi\,\eh .
\end{align}
All three are the cross product of one angular velocity with the basis vector,
$d\vb b/dt = \vb\omega_{LE}\times\vb b$, where
\begin{equation}
  \vb\omega_{LE} = \dot\theta\sin\phi\,\uh \;-\; \dot\phi\,\eh \;+\; \dot\theta\cos\phi\,\nh
  \label{eq:omegaLE}
\end{equation}
is the angular velocity of the local frame relative to Earth.

\section{Translational kinematics}

Differentiating $\vb r = r\uh$ in the Earth frame and using \cref{eq:duh}:
\begin{equation}
  \vb v = \dot r\,\uh + r\frac{d\uh}{dt}
        = \dot r\,\uh + r\dot\phi\,\nh + r\dot\theta\cos\phi\,\eh .
  \label{eq:v_components}
\end{equation}
The speed $V$, flight-path angle $\gamma$ and heading $\psi$ are \emph{defined} by writing the
same velocity as
\begin{equation}
  \vb v = V\tth,\qquad
  \tth = \sin\gamma\,\uh + \cos\gamma\cos\psi\,\nh + \cos\gamma\sin\psi\,\eh .
  \label{eq:that}
\end{equation}
Matching the $\uh$, $\nh$ and $\eh$ components of \cref{eq:v_components,eq:that}:
\begin{align}
  \dot r &= V\sin\gamma, \label{eq:rdot}\\
  \dot\phi &= \frac{V\cos\gamma\cos\psi}{r}, \label{eq:phidot}\\
  \dot\theta &= \frac{V\cos\gamma\sin\psi}{r\cos\phi}. \label{eq:thetadot}
\end{align}

\section{Velocity dynamics}\label{sec:veldyn}

Two more unit vectors complete an orthonormal triad with $\tth$. They are the
directions in which $\tth$ moves when $\gamma$ or $\psi$ changes:
\begin{align}
  \gh &\equiv \pd{\tth}{\gamma} = \cos\gamma\,\uh - \sin\gamma\cos\psi\,\nh - \sin\gamma\sin\psi\,\eh,\\
  \hh &\equiv \frac{1}{\cos\gamma}\pd{\tth}{\psi} = -\sin\psi\,\nh + \cos\psi\,\eh .
\end{align}
$\gh$ lies in the vertical plane, perpendicular to the velocity on the upper side; $\hh$
is horizontal and to the right of the velocity. One can check that
$|\gh|=|\hh|=1$ and $\tth\cdot\gh=\tth\cdot\hh=\gh\cdot\hh=0$.

Differentiate $\vb v = V\tth$ in the Earth frame. The unit vector $\tth$ changes for two
reasons: its components change because $\gamma$ and $\psi$ change, and the local basis it
is written in turns with $\vb\omega_{LE}$. So
\begin{align}
  \left.\frac{d\tth}{dt}\right|_E &= \dot\gamma\,\pd{\tth}{\gamma} + \dot\psi\,\pd{\tth}{\psi}
     + \vb\omega_{LE}\times\tth
   = \dot\gamma\,\gh + \dot\psi\cos\gamma\,\hh + \vb\omega_{LE}\times\tth, \\[4pt]
  \left.\frac{d\vb v}{dt}\right|_E &= \dot V\,\tth + V\dot\gamma\,\gh + V\dot\psi\cos\gamma\,\hh
     + V\,\vb\omega_{LE}\times\tth . \label{eq:dvdt}
\end{align}

\paragraph{Transport terms.} Substituting \cref{eq:phidot,eq:thetadot} into
\cref{eq:omegaLE} and projecting $\vb\omega_{LE}\times\tth$ onto the triad gives
\begin{align}
  \tth\cdot(\vb\omega_{LE}\times\tth) &= 0,\\
  \gh\cdot(\vb\omega_{LE}\times\tth) &= -\frac{V}{r}\cos\gamma,\\
  \hh\cdot(\vb\omega_{LE}\times\tth) &= -\frac{V}{r}\cos^2\!\gamma\,\sin\psi\tan\phi .
\end{align}
These are the effects of flying over a curved Earth: the local horizon tilts as the vehicle
moves, and the direction of north changes with longitude.

\paragraph{Projection.} Let $\vb a$ be the right-hand side of \cref{eq:rotating}. Dotting
\cref{eq:dvdt} with $\tth$, $\gh$ and $\hh$ in turn, and moving the transport terms to the
right:
\begin{align}
  \dot V &= \tth\cdot\vb a, \label{eq:Vdot_proj}\\
  \dot\gamma &= \frac{1}{V}\,\gh\cdot\vb a + \frac{V}{r}\cos\gamma, \label{eq:gdot_proj}\\
  \dot\psi &= \frac{1}{V\cos\gamma}\,\hh\cdot\vb a + \frac{V}{r}\cos\gamma\sin\psi\tan\phi .
  \label{eq:pdot_proj}
\end{align}
What remains is to project each force term of $\vb a$.

\section{Forces}

\subsection{Aerodynamic force}

Drag acts opposite the velocity. Lift is perpendicular to it, rotated from $\gh$ toward
$\hh$ by the bank angle:
\begin{equation}
  \vb a_{\text{aero}} = -\aD\,\tth + \aL\big(\cos\sigma\,\gh + \sin\sigma\,\hh\big),
\end{equation}
\begin{equation}
  \aL = \frac{\rho(h)\,V^2 S}{2m}\,C_L(\alpha, M),\qquad
  \aD = \frac{\rho(h)\,V^2 S}{2m}\,C_D(\alpha, M),\qquad
  M = \frac{V}{a_s(h)} .
  \label{eq:aero}
\end{equation}
Projections:
\[
  \tth\cdot\vb a_{\text{aero}} = -\aD,\qquad
  \gh\cdot\vb a_{\text{aero}} = \aL\cos\sigma,\qquad
  \hh\cdot\vb a_{\text{aero}} = \aL\sin\sigma .
\]
With $\sigma>0$ the lift has a component along $+\hh$ (to the right), which turns the
vehicle right.

\subsection{Gravity with $J_2$}

The gravitational potential truncated after $J_2$ is
\begin{equation}
  U(r,\phi) = \frac{\mu}{r}\left[1 - J_2\left(\frac{R_E}{r}\right)^2 P_2(\sin\phi)\right],
  \qquad P_2(x) = \tfrac12\left(3x^2-1\right),
\end{equation}
and $\vb g = \nabla U$. In the local frame, the radial and northward components are
\begin{align}
  \uh\cdot\vb g &= \pd{U}{r}
   = -\frac{\mu}{r^2} + \frac{3\mu J_2 R_E^2}{r^4}P_2(\sin\phi)
   = -\frac{\mu}{r^2}\left[1 - \frac32 J_2\left(\frac{R_E}{r}\right)^2\left(3\sin^2\phi-1\right)\right],\\
  \nh\cdot\vb g &= \frac1r\pd{U}{\phi}
   = \frac{\mu}{r^2}\left(-J_2\frac{R_E^2}{r^2}\right)\frac{d P_2}{d\phi}
   = -\frac{3\mu J_2 R_E^2}{r^4}\sin\phi\cos\phi .
\end{align}
Define the downward radial magnitude $g_r$ and the signed northward component $g_n$:
\begin{equation}
  g_r = \frac{\mu}{r^2}\left[1 - \frac32 J_2\left(\frac{R_E}{r}\right)^2\left(3\sin^2\phi-1\right)\right],
  \qquad
  g_n = -\frac{3\mu J_2 R_E^2}{r^4}\sin\phi\cos\phi,
  \qquad
  \vb g = -g_r\,\uh + g_n\,\nh .
  \label{eq:gravity}
\end{equation}
Projections:
\[
  \tth\cdot\vb g = -g_r\sin\gamma + g_n\cos\gamma\cos\psi,\qquad
  \gh\cdot\vb g = -g_r\cos\gamma - g_n\sin\gamma\cos\psi,\qquad
  \hh\cdot\vb g = -g_n\sin\psi .
\]

\subsection{Coriolis and centripetal terms}

Using \cref{eq:omega_local} for $\vb\Omega$, $\vb v = V\tth$ and $\vb r = r\uh$, the
Coriolis acceleration $-2\vb\Omega\times\vb v$ projects as
\begin{align}
  \tth\cdot\left(-2\vb\Omega\times\vb v\right) &= 0, \\
  \gh\cdot\left(-2\vb\Omega\times\vb v\right) &= 2\Omega V\cos\phi\sin\psi, \\
  \hh\cdot\left(-2\vb\Omega\times\vb v\right) &= 2\Omega V\left(\sin\phi\cos\gamma - \cos\phi\sin\gamma\cos\psi\right),
\end{align}
and the centripetal acceleration $-\vb\Omega\times(\vb\Omega\times\vb r)$ as
\begin{align}
  \tth\cdot\left(-\vb\Omega\times(\vb\Omega\times\vb r)\right)
    &= \Omega^2 r\cos\phi\left(\sin\gamma\cos\phi - \cos\gamma\sin\phi\cos\psi\right), \\
  \gh\cdot\left(-\vb\Omega\times(\vb\Omega\times\vb r)\right)
    &= \Omega^2 r\cos\phi\left(\cos\gamma\cos\phi + \sin\gamma\sin\phi\cos\psi\right), \\
  \hh\cdot\left(-\vb\Omega\times(\vb\Omega\times\vb r)\right)
    &= \Omega^2 r\sin\phi\cos\phi\sin\psi .
\end{align}
The Coriolis term has no component along the velocity, so it never changes the speed.

\section{The assembled nonlinear model}\label{sec:full_eom}

Substituting all projections into \cref{eq:Vdot_proj,eq:gdot_proj,eq:pdot_proj} and
collecting the kinematics \cref{eq:rdot,eq:phidot,eq:thetadot}:
\begin{align}
  \dot r &= V\sin\gamma \label{eq:full_r}\\[4pt]
  \dot\theta &= \frac{V\cos\gamma\sin\psi}{r\cos\phi} \\[4pt]
  \dot\phi &= \frac{V\cos\gamma\cos\psi}{r} \\[4pt]
  \dot V &= -\aD - g_r\sin\gamma + g_n\cos\gamma\cos\psi
            + \Omega^2 r\cos\phi\left(\sin\gamma\cos\phi - \cos\gamma\sin\phi\cos\psi\right) \\[4pt]
  \dot\gamma &= \frac{\aL\cos\sigma}{V} - \frac{g_r\cos\gamma}{V} - \frac{g_n\sin\gamma\cos\psi}{V}
            + \frac{V\cos\gamma}{r} + 2\Omega\cos\phi\sin\psi \notag\\
         &\qquad + \frac{\Omega^2 r\cos\phi}{V}\left(\cos\gamma\cos\phi + \sin\gamma\sin\phi\cos\psi\right) \\[4pt]
  \dot\psi &= \frac{\aL\sin\sigma}{V\cos\gamma} - \frac{g_n\sin\psi}{V\cos\gamma}
            + \frac{V\cos\gamma\sin\psi\tan\phi}{r}
            - 2\Omega\left(\cos\phi\cos\psi\tan\gamma - \sin\phi\right) \notag\\
         &\qquad + \frac{\Omega^2 r\sin\phi\cos\phi\sin\psi}{V\cos\gamma} \label{eq:full_psi}
\end{align}
In the planner, two more equations make $\alpha$ and $\sigma$ states:
\begin{equation}
  \dot\alpha = u_\alpha,\qquad \dot\sigma = u_\sigma .
\end{equation}

These are line for line the expressions in \code{HL20Model.xdot}.

\section{Aerodynamic and atmospheric sub-models}

\begin{itemize}
\item \textbf{Altitude.} The atmosphere is a function of geodetic altitude. The model uses
  the height above the WGS84 ellipsoid along the radius,
  \begin{equation}
    h = r - R_{\text{ell}}(\phi),\qquad
    R_{\text{ell}}(\phi) = \frac{a_W b_W}{\sqrt{b_W^2\cos^2\phi + a_W^2\sin^2\phi}},\qquad
    b_W = a_W(1-f_W),
  \end{equation}
  which differs from the true geodetic altitude by much less than a meter here. Using
  $r - R_E$ instead would be wrong by about 5 km at $30^\circ$ latitude, which is a
  factor of about 2 in density.
\item \textbf{Atmosphere.} $\ln\rho(h)$ and $a_s(h)$ are B-spline fits of
  \code{hgv3.atmosphere} (USSA76) sampled every 500 m (\code{atmosphereFits}).
\item \textbf{Aerodynamics.} $C_L(\alpha,M)$ and $C_D(\alpha,M)$ are 2-D B-spline fits of
  \code{hgv3.wind\_coefficients} on an $(\alpha,M)$ grid (\code{aeroFits}). The splines give
  the continuous first and second derivatives the NLP solver needs; the simulator's own
  table is linear in Mach with clamped $\alpha$.
\item \textbf{Valid $\alpha$ window.} \code{hgv3} clamps $\alpha$ to each Mach row's range
  and blends rows linearly in Mach. The window in which no clamping occurs is the
  intersection of the two bracketing rows' ranges; \code{alphaLimitFits} builds a
  continuous, conservative version $\alpha_{\text{lo}}(M)\le\alpha\le\alpha_{\text{hi}}(M)$.
\end{itemize}

\section{Relation to the \texttt{hgv3} simulator}

\code{hgv3} integrates \cref{eq:newton} in ECI with fixed-step RK4. The physics is the same
(same aero table, USSA76, point-mass plus $J_2$ gravity), with two implementation
differences that matter:
\begin{enumerate}
\item It co-rotates the atmosphere about the ECI $z$ axis, while its ECEF frame comes from
  the full ERFA precession--nutation chain. In 2026 the two axes differ by $0.148^\circ$,
  so the simulated air mass moves about 1.2 m/s relative to ECEF at the equator. The
  model above co-rotates about the ECEF pole.
\item Its aero coefficients are piecewise linear in Mach with $\alpha$ clamping, where the
  model uses smooth fits.
\end{enumerate}
Flying the same open-loop $(\alpha,\sigma)$ for 800 s from Mach 15, the two agree to about
10 m in altitude and 0.1 m/s in speed, with an along-track drift of about 0.8 m/s
(670 m after 800 s) consistent with item 1.

%======================================================================
\chapter{Linearization about a planned trajectory}\label{ch:lin}

\section{Reference trajectory and perturbations}

For guidance, write the model as $\dot{\vb x} = \vb f(\vb x,\vb u)$ with
\[
  \vb x = [\,r,\ \theta,\ \phi,\ V,\ \gamma,\ \psi\,]^T,\qquad
  \vb u = [\,\alpha,\ \sigma\,]^T .
\]
The planner supplies a reference $\vb x^*(t)$, $\vb u^*(t)$ that satisfies the model:
$\dot{\vb x}^* = \vb f(\vb x^*,\vb u^*)$. Define the perturbations
\[
  \delta\vb x = \vb x - \vb x^*(t),\qquad \delta\vb u = \vb u - \vb u^*(t).
\]
A first-order Taylor expansion of $\vb f$ about the reference gives
\begin{equation}
  \delta\dot{\vb x} = A(t)\,\delta\vb x + B(t)\,\delta\vb u,\qquad
  A(t) = \left.\pd{\vb f}{\vb x}\right|_{*},\qquad
  B(t) = \left.\pd{\vb f}{\vb u}\right|_{*},
\end{equation}
where $|_*$ means evaluated on the reference. The zeroth-order terms cancel precisely because
the reference satisfies the model.

\section{Simplifications for the hand derivation}\label{sec:simplify}

For the hand derivation only, Earth rotation and $J_2$ are dropped, and altitude is taken
as $h = r - R_0$ with $R_0$ the local ellipsoid radius held constant:
\begin{equation}
  \Omega = 0,\qquad J_2 = 0\ \Rightarrow\ g_r = g \equiv \frac{\mu}{r^2},\ g_n = 0,\qquad
  \pd{h}{r} = 1,\ \pd{h}{\phi} = 0 .
\end{equation}
This was checked on the test plan by computing exact Jacobians of the full model and of the
simplified one with CasADi:
\begin{itemize}
\item With states scaled to typical error sizes (1 km of position, 10 m/s, $1^\circ$), the
  largest entry of $A$ is about $0.11$ s$^{-1}$, and rotation plus $J_2$ change any entry by at
  most $5\times10^{-5}$ s$^{-1}$. $B$ is unchanged.
\item The open-loop eigenvalues move by at most $2\times10^{-3}$ s$^{-1}$.
\item The neglected latitude dependence of $h$ is about 3 m of altitude per km of northward
  displacement.
\end{itemize}
Rotation and $J_2$ do dominate a few entries that are themselves tiny (for example
$\partial\dot\gamma/\partial\psi$ comes only from the Coriolis term), but those entries
correspond to time scales of hours. The simplified model therefore captures the structure
that matters for control. In code, use the exact Jacobians (\cref{sec:casadi}).

The simplified equations are:
\begin{align}
  \dot r &= V\sin\gamma \label{eq:s_r}\\
  \dot\theta &= \frac{V\cos\gamma\sin\psi}{r\cos\phi} \\
  \dot\phi &= \frac{V\cos\gamma\cos\psi}{r} \\
  \dot V &= -\aD - g\sin\gamma \label{eq:s_V}\\
  \dot\gamma &= \frac{\aL\cos\sigma}{V} - \frac{g\cos\gamma}{V} + \frac{V\cos\gamma}{r} \label{eq:s_g}\\
  \dot\psi &= \frac{\aL\sin\sigma}{V\cos\gamma} + \frac{V\cos\gamma\sin\psi\tan\phi}{r} \label{eq:s_p}
\end{align}

\section{Aerodynamic partial derivatives}

Write $k \equiv \bar q S/m = \rho V^2 S/(2m)$, so that $\aL = k\,C_L(\alpha,M)$ and
$\aD = k\,C_D(\alpha,M)$. The coefficient slopes and atmospheric gradients are
\[
  C_{L,\alpha} = \pd{C_L}{\alpha},\quad C_{L,M} = \pd{C_L}{M},\quad
  C_{D,\alpha} = \pd{C_D}{\alpha},\quad C_{D,M} = \pd{C_D}{M},
\]
\[
  H_\rho \equiv -\frac{\rho}{d\rho/dh}\quad\text{(local density scale height)},\qquad
  a_s' \equiv \frac{d a_s}{dh} .
\]
Partials of the intermediate quantities:
\begin{align}
  \pd{k}{V} &= \frac{\rho V S}{m} = \frac{2k}{V}, &
  \pd{k}{h} &= \frac{d\rho}{dh}\frac{V^2S}{2m} = -\frac{k}{H_\rho}, \\
  \pd{M}{V} &= \frac{1}{a_s}, &
  \pd{M}{h} &= -\frac{V a_s'}{a_s^2} .
\end{align}
By the chain rule, with $\aL = k\,C_L(\alpha, M(V,h))$:
\begin{align}
  \pd{\aL}{V} &= \pd{k}{V}C_L + k\,C_{L,M}\pd{M}{V}
              = \frac{2\aL}{V} + \frac{k\,C_{L,M}}{a_s}, \label{eq:dLdV}\\
  \pd{\aL}{h} &= \pd{k}{h}C_L + k\,C_{L,M}\pd{M}{h}
              = -\frac{\aL}{H_\rho} - \frac{k\,C_{L,M}\,V a_s'}{a_s^2}, \label{eq:dLdh}\\
  \pd{\aL}{\alpha} &= k\,C_{L,\alpha}, \label{eq:dLda}
\end{align}
and identically for drag:
\begin{align}
  \pd{\aD}{V} &= \frac{2\aD}{V} + \frac{k\,C_{D,M}}{a_s}, &
  \pd{\aD}{h} &= -\frac{\aD}{H_\rho} - \frac{k\,C_{D,M}\,V a_s'}{a_s^2}, &
  \pd{\aD}{\alpha} &= k\,C_{D,\alpha} .
\end{align}
Gravity: $\partial g/\partial r = -2g/r$. Since $h = r - R_0$, a partial with respect to $r$
equals the partial with respect to $h$.

\section{Partial derivatives of the kinematic rows}

\paragraph{$\dot r = V\sin\gamma$:}
\[
  \pd{\dot r}{V} = \sin\gamma,\qquad \pd{\dot r}{\gamma} = V\cos\gamma .
\]
\paragraph{$\dot\theta = V\cos\gamma\sin\psi/(r\cos\phi)$:}
\begin{align*}
  \pd{\dot\theta}{r} &= -\frac{V\cos\gamma\sin\psi}{r^2\cos\phi}, &
  \pd{\dot\theta}{\phi} &= \frac{V\cos\gamma\sin\psi\tan\phi}{r\cos\phi}, &
  \pd{\dot\theta}{V} &= \frac{\cos\gamma\sin\psi}{r\cos\phi}, \\
  \pd{\dot\theta}{\gamma} &= -\frac{V\sin\gamma\sin\psi}{r\cos\phi}, &
  \pd{\dot\theta}{\psi} &= \frac{V\cos\gamma\cos\psi}{r\cos\phi} .
\end{align*}
\paragraph{$\dot\phi = V\cos\gamma\cos\psi/r$:}
\begin{align*}
  \pd{\dot\phi}{r} &= -\frac{V\cos\gamma\cos\psi}{r^2}, &
  \pd{\dot\phi}{V} &= \frac{\cos\gamma\cos\psi}{r}, &
  \pd{\dot\phi}{\gamma} &= -\frac{V\sin\gamma\cos\psi}{r}, &
  \pd{\dot\phi}{\psi} &= -\frac{V\cos\gamma\sin\psi}{r} .
\end{align*}
No kinematic row depends on $\theta$, $\alpha$ or $\sigma$.

\section{Partial derivatives of the dynamic rows}

Each partial is named $a_{ij}$ (state) or $b_{ij}$ (input) for use in \cref{ch:err}.

\paragraph{Speed, $\dot V = -\aD - g\sin\gamma$:}
\begin{align}
  a_{Vr} &= \pd{\dot V}{r} = -\pd{\aD}{h} + \frac{2g}{r}\sin\gamma, &
  a_{VV} &= \pd{\dot V}{V} = -\pd{\aD}{V}, &
  a_{V\gamma} &= \pd{\dot V}{\gamma} = -g\cos\gamma, \\
  b_{V\alpha} &= \pd{\dot V}{\alpha} = -\pd{\aD}{\alpha}, &
  b_{V\sigma} &= \pd{\dot V}{\sigma} = 0 . \notag
\end{align}

\paragraph{Flight-path angle, $\dot\gamma = \aL\cos\sigma/V - g\cos\gamma/V + V\cos\gamma/r$:}
\begin{align}
  a_{\gamma r} &= \pd{\dot\gamma}{r} = \frac{\cos\sigma}{V}\pd{\aL}{h}
     + \left(\frac{2g}{rV} - \frac{V}{r^2}\right)\cos\gamma, \\
  a_{\gamma V} &= \pd{\dot\gamma}{V} = \frac{\cos\sigma}{V}\pd{\aL}{V} - \frac{\aL\cos\sigma}{V^2}
     + \left(\frac{g}{V^2} + \frac1r\right)\cos\gamma, \\
  a_{\gamma\gamma} &= \pd{\dot\gamma}{\gamma} = \left(\frac{g}{V} - \frac{V}{r}\right)\sin\gamma, \\
  b_{\gamma\alpha} &= \pd{\dot\gamma}{\alpha} = \frac{\cos\sigma}{V}\pd{\aL}{\alpha}, \qquad
  b_{\gamma\sigma} = \pd{\dot\gamma}{\sigma} = -\frac{\aL\sin\sigma}{V} .
\end{align}

\paragraph{Heading, $\dot\psi = \aL\sin\sigma/(V\cos\gamma) + V\cos\gamma\sin\psi\tan\phi/r$:}
\begin{align}
  a_{\psi r} &= \pd{\dot\psi}{r} = \frac{\sin\sigma}{V\cos\gamma}\pd{\aL}{h}
     - \frac{V}{r^2}\cos\gamma\sin\psi\tan\phi, \\
  a_{\psi\phi} &= \pd{\dot\psi}{\phi} = \frac{V}{r}\frac{\cos\gamma\sin\psi}{\cos^2\phi}, \\
  a_{\psi V} &= \pd{\dot\psi}{V} = \frac{\sin\sigma}{V\cos\gamma}\pd{\aL}{V}
     - \frac{\aL\sin\sigma}{V^2\cos\gamma} + \frac{1}{r}\cos\gamma\sin\psi\tan\phi, \\
  a_{\psi\gamma} &= \pd{\dot\psi}{\gamma} = \frac{\aL\sin\sigma\sin\gamma}{V\cos^2\gamma}
     - \frac{V}{r}\sin\gamma\sin\psi\tan\phi, \\
  a_{\psi\psi} &= \pd{\dot\psi}{\psi} = \frac{V}{r}\cos\gamma\cos\psi\tan\phi, \\
  b_{\psi\alpha} &= \pd{\dot\psi}{\alpha} = \frac{\sin\sigma}{V\cos\gamma}\pd{\aL}{\alpha}, \qquad
  b_{\psi\sigma} = \pd{\dot\psi}{\sigma} = \frac{\aL\cos\sigma}{V\cos\gamma} .
\end{align}
The speed and flight-path-angle rows do not depend on $\theta$, $\phi$ or $\psi$ in the
simplified model.

\section{The linear model}

With the state order $[r,\theta,\phi,V,\gamma,\psi]$ and input order $[\alpha,\sigma]$,
writing the kinematic partials as $a_{r V}=\partial\dot r/\partial V$, and so on:
\begin{equation}
\begin{bmatrix}\delta\dot r\\ \delta\dot\theta\\ \delta\dot\phi\\ \delta\dot V\\ \delta\dot\gamma\\ \delta\dot\psi\end{bmatrix}
=
\begin{bmatrix}
0 & 0 & 0 & a_{rV} & a_{r\gamma} & 0\\
a_{\theta r} & 0 & a_{\theta\phi} & a_{\theta V} & a_{\theta\gamma} & a_{\theta\psi}\\
a_{\phi r} & 0 & 0 & a_{\phi V} & a_{\phi\gamma} & a_{\phi\psi}\\
a_{Vr} & 0 & 0 & a_{VV} & a_{V\gamma} & 0\\
a_{\gamma r} & 0 & 0 & a_{\gamma V} & a_{\gamma\gamma} & 0\\
a_{\psi r} & 0 & a_{\psi\phi} & a_{\psi V} & a_{\psi\gamma} & a_{\psi\psi}
\end{bmatrix}
\begin{bmatrix}\delta r\\ \delta\theta\\ \delta\phi\\ \delta V\\ \delta\gamma\\ \delta\psi\end{bmatrix}
+
\begin{bmatrix}
0 & 0\\ 0 & 0\\ 0 & 0\\ b_{V\alpha} & 0\\ b_{\gamma\alpha} & b_{\gamma\sigma}\\ b_{\psi\alpha} & b_{\psi\sigma}
\end{bmatrix}
\begin{bmatrix}\delta\alpha\\ \delta\sigma\end{bmatrix}.
\label{eq:linmodel}
\end{equation}

\section{Structure of the linear model}

\begin{itemize}
\item \textbf{Longitude is ignorable.} The $\theta$ column is zero: nothing depends on
  longitude, so it is a pure integrator (an eigenvalue at zero). This is why tracking
  in latitude and longitude is awkward, and why \cref{ch:err} uses errors measured along
  and across the path instead.
\item \textbf{Longitudinal block.} $(\delta r, \delta V, \delta\gamma)$ form a closed
  subsystem driven by $\delta\alpha$ (through $b_{V\alpha}$ and $b_{\gamma\alpha}$) and by
  $\delta\sigma$ (through $b_{\gamma\sigma}$).
\item \textbf{Lateral block.} $(\delta\theta, \delta\phi, \delta\psi)$ are driven by the
  longitudinal states and by $\delta\sigma$ through $b_{\psi\sigma}$.
\item \textbf{Coupling through bank.} When $\sigma^*\approx 0$: $b_{\gamma\sigma}\approx 0$
  and $b_{\psi\alpha}\approx 0$, so $\alpha$ controls the longitudinal motion and $\sigma$ the
  lateral motion, separately. This is the analog of Rajamani's decoupled lateral and
  longitudinal vehicle models. With large $\sigma^*$, a change in bank shifts lift between
  vertical and horizontal, which couples the two. Entry guidance exploits exactly this:
  bank \emph{magnitude} controls range, and bank \emph{sign} controls crossrange.
\item \textbf{Open-loop modes.} At one point on the test plan (Mach about 11), the
  eigenvalues are a lightly damped oscillation in speed, altitude and flight-path angle
  at about $0.026$ rad/s (period near 4 min, damping ratio about 0.06; the entry analog of
  the phugoid), slow modes near $10^{-3}$ s$^{-1}$ (one slightly unstable), and the zero
  eigenvalue from longitude.
\end{itemize}

%======================================================================
\chapter{Error-state model for path following}\label{ch:err}

This chapter goes back to the nonlinear equations of \cref{sec:simplify}, substitutes
errors measured relative to the reference \emph{path}, and linearizes. It mirrors the
derivation of Rajamani's error-state lateral model ($e_1$, $e_2$), with the speed loop
folded into the same model because a glider's only longitudinal lever is drag.

\section{Reference path}

From the planned trajectory, compute the ground-track arc length $s$ along a sphere of
radius $R_0$. Along the reference,
\begin{equation}
  \sdot^* = \frac{R_0}{r^*}\,V^*\cos\gamma^* ,
\end{equation}
and every reference quantity can be expressed as a function of $s$:
$h^*(s)$, $V^*(s)$, $\gamma^*(s)$, $\psi^*(s)$, $\alpha^*(s)$, $\sigma^*(s)$. Define also the
reference rates, which come straight from the model evaluated on the reference,
\[
  \dot V^* = f_V(\vb x^*,\vb u^*),\qquad
  \dot\gamma^* = f_\gamma(\vb x^*,\vb u^*),\qquad
  \dot\psi^* = f_\psi(\vb x^*,\vb u^*),
\]
and the geodesic curvature of the ground track, positive for a right turn:
\begin{equation}
  \kappa(s) = \frac{d\psi^*}{ds} - \frac{\sin\psi^*\tan\phi^*}{R_0}.
\end{equation}
The second term removes the change in heading caused only by the convergence of meridians;
$\kappa$ is the curvature a local flat map would show.

\paragraph{Reference point.} At each instant, $s(t)$ is chosen as the point on the reference
ground track closest to the vehicle, so the along-track position error is zero by
construction. How that point is found (projection, spline search, integrating $\sdot$) is
independent of what follows.

\section{Error states}

\begin{center}
\begin{tabular}{@{}lll@{}}
\toprule
Error & Definition & Meaning \\
\midrule
$e_c$ & signed distance from the reference ground track & crossrange error, positive to the right \\
$e_h$ & $h - h^*(s)$ & altitude error \\
$e_\psi$ & $\psi - \psi^*(s)$ & heading error \\
$e_\gamma$ & $\gamma - \gamma^*(s)$ & flight-path angle error \\
$e_V$ & $V - V^*(s)$ & speed error at the same path position \\
\midrule
$\delta\alpha$ & $\alpha - \alpha^*(s)$ & angle-of-attack correction \\
$\delta\sigma$ & $\sigma - \sigma^*(s)$ & bank correction \\
\bottomrule
\end{tabular}
\end{center}
The vehicle state is therefore $V = V^*+e_V$, $\gamma = \gamma^*+e_\gamma$,
$\psi=\psi^*+e_\psi$, $r = r^*+e_h$, $\alpha = \alpha^*+\delta\alpha$, $\sigma=\sigma^*+\delta\sigma$,
all reference values taken at the current $s$.

\section{Rate of progress along the path}

The vehicle's horizontal velocity, scaled down to the reference sphere, has magnitude
$(R_0/r)V\cos\gamma$. Its component along the path direction is that times $\cos e_\psi$.
A curve offset by $e_c$ toward the center of a turn of curvature $\kappa$ is shorter by the
factor $(1-\kappa e_c)$, so the projected point advances faster by the inverse factor:
\begin{equation}
  \sdot = \frac{R_0}{r}\,\frac{V\cos\gamma\cos e_\psi}{1 - \kappa e_c} .
  \label{eq:sdot}
\end{equation}
The component across the path changes the crossrange error:
\begin{equation}
  \dot e_c = \frac{R_0}{r}\,V\cos\gamma\sin e_\psi .
  \label{eq:ecdot}
\end{equation}

\section{Nonlinear error dynamics}

Differentiate each error. A reference quantity $X^*(s)$ changes only because $s$ changes:
$\frac{d}{dt}X^*(s) = \frac{dX^*}{ds}\,\sdot$. So
\begin{align}
  \dot e_c &= \frac{R_0}{r}\,V\cos\gamma\sin e_\psi, \label{eq:ne_c}\\
  \dot e_h &= V\sin\gamma - \frac{dh^*}{ds}\,\sdot, \label{eq:ne_h}\\
  \dot e_\psi &= f_\psi(\vb x,\vb u) - \frac{d\psi^*}{ds}\,\sdot, \\
  \dot e_\gamma &= f_\gamma(\vb x,\vb u) - \frac{d\gamma^*}{ds}\,\sdot, \\
  \dot e_V &= f_V(\vb x,\vb u) - \frac{dV^*}{ds}\,\sdot, \label{eq:ne_V}
\end{align}
with $\sdot$ from \cref{eq:sdot} and $f_V$, $f_\gamma$, $f_\psi$ the right-hand sides of
\cref{eq:s_V,eq:s_g,eq:s_p} evaluated at the \emph{actual} state and input. The reference
slopes follow from dividing reference rates by $\sdot^*$:
\begin{equation}
  \frac{dh^*}{ds} = \frac{V^*\sin\gamma^*}{\sdot^*} = \frac{r^*}{R_0}\tan\gamma^*,\qquad
  \frac{d\psi^*}{ds} = \frac{\dot\psi^*}{\sdot^*},\qquad
  \frac{d\gamma^*}{ds} = \frac{\dot\gamma^*}{\sdot^*},\qquad
  \frac{dV^*}{ds} = \frac{\dot V^*}{\sdot^*} .
\end{equation}
These play the role of Rajamani's $\dot\psi_{\text{des}}$ terms: they are the rates at which
the reference itself moves.

\section{Linearization}

\subsection*{Step 1: linearize the path rate}

Expand \cref{eq:sdot} to first order about zero error. Each partial is evaluated at the
reference:
\begin{align}
  \pd{\sdot}{e_V} &= \frac{R_0}{r^*}\cos\gamma^* = \frac{\sdot^*}{V^*}, &
  \pd{\sdot}{e_\gamma} &= -\frac{R_0}{r^*}V^*\sin\gamma^* = -\sdot^*\tan\gamma^*, \notag\\
  \pd{\sdot}{e_h} &= -\frac{R_0}{r^{*2}}V^*\cos\gamma^* = -\frac{\sdot^*}{r^*}, &
  \pd{\sdot}{e_c} &= \sdot^*\kappa, \qquad
  \pd{\sdot}{e_\psi} = 0 . \notag
\end{align}
The heading-error partial vanishes because $\partial\cos e_\psi/\partial e_\psi = 0$ at
$e_\psi=0$. So
\begin{equation}
  \sdot \approx \sdot^*\left(1 + \frac{e_V}{V^*} - \tan\gamma^*\,e_\gamma - \frac{e_h}{r^*} + \kappa\,e_c\right).
  \label{eq:sdot_lin}
\end{equation}

\subsection*{Step 2: crossrange row}

From \cref{eq:ne_c}, with $\sin e_\psi\approx e_\psi$:
\begin{equation}
  \dot e_c = \sdot^*\,e_\psi .
\end{equation}

\subsection*{Step 3: altitude row}

Linearize \cref{eq:ne_h} term by term. The first term gives
$\delta(V\sin\gamma) = \sin\gamma^*\,e_V + V^*\cos\gamma^*\,e_\gamma$. The second term is
$\frac{r^*}{R_0}\tan\gamma^*$ times \cref{eq:sdot_lin}. Collecting by error:
\begin{align*}
  e_V:&\quad \sin\gamma^* - \frac{r^*}{R_0}\tan\gamma^*\,\frac{\sdot^*}{V^*}
        = \sin\gamma^* - \sin\gamma^* = 0,\\
  e_\gamma:&\quad V^*\cos\gamma^* + \frac{r^*}{R_0}\tan\gamma^*\,\sdot^*\tan\gamma^*
        = V^*\cos\gamma^* + V^*\frac{\sin^2\gamma^*}{\cos\gamma^*} = \frac{V^*}{\cos\gamma^*},\\
  e_h:&\quad \frac{r^*}{R_0}\tan\gamma^*\,\frac{\sdot^*}{r^*} = \frac{V^*\sin\gamma^*}{r^*},\\
  e_c:&\quad -\frac{r^*}{R_0}\tan\gamma^*\,\sdot^*\kappa = -V^*\kappa\sin\gamma^* .
\end{align*}
Speed error drops out: going faster along the same path changes where you are along it,
not your height above it. So
\begin{equation}
  \dot e_h = -V^*\kappa\sin\gamma^*\,e_c + \frac{V^*\sin\gamma^*}{r^*}\,e_h + \frac{V^*}{\cos\gamma^*}\,e_\gamma .
\end{equation}

\subsection*{Step 4: heading, flight-path angle and speed rows}

For each of the three dynamic rows, $f_\square(\vb x,\vb u)$ linearizes to the
\cref{ch:lin} partials, and the path-rate term contributes $-\dot\square^*/\sdot^*$ times
\cref{eq:sdot_lin}. For heading:
\begin{align}
  \dot e_\psi ={}& a_{\psi r}\,e_h + a_{\psi V}\,e_V + a_{\psi\gamma}\,e_\gamma + a_{\psi\psi}\,e_\psi
       + b_{\psi\alpha}\,\delta\alpha + b_{\psi\sigma}\,\delta\sigma \notag\\
     & - \dot\psi^*\left(\frac{e_V}{V^*} - \tan\gamma^*\,e_\gamma - \frac{e_h}{r^*} + \kappa\,e_c\right),
\end{align}
and collecting terms:
\begin{align}
  \dot e_\psi ={}& -\kappa\,\dot\psi^*\,e_c
    + \left(a_{\psi r} + \frac{\dot\psi^*}{r^*}\right)e_h
    + a_{\psi\psi}\,e_\psi
    + \left(a_{\psi\gamma} + \dot\psi^*\tan\gamma^*\right)e_\gamma
    + \left(a_{\psi V} - \frac{\dot\psi^*}{V^*}\right)e_V \notag\\
   & + b_{\psi\alpha}\,\delta\alpha + b_{\psi\sigma}\,\delta\sigma, \\[6pt]
  \dot e_\gamma ={}& -\kappa\,\dot\gamma^*\,e_c
    + \left(a_{\gamma r} + \frac{\dot\gamma^*}{r^*}\right)e_h
    + \left(a_{\gamma\gamma} + \dot\gamma^*\tan\gamma^*\right)e_\gamma
    + \left(a_{\gamma V} - \frac{\dot\gamma^*}{V^*}\right)e_V \notag\\
   & + b_{\gamma\alpha}\,\delta\alpha + b_{\gamma\sigma}\,\delta\sigma, \\[6pt]
  \dot e_V ={}& -\kappa\,\dot V^*\,e_c
    + \left(a_{Vr} + \frac{\dot V^*}{r^*}\right)e_h
    + \left(a_{V\gamma} + \dot V^*\tan\gamma^*\right)e_\gamma
    + \left(a_{VV} - \frac{\dot V^*}{V^*}\right)e_V
    + b_{V\alpha}\,\delta\alpha .
\end{align}
All $a_{ij}$ and $b_{ij}$ are the \cref{ch:lin} expressions evaluated on the reference at
the current $s$. The heading error does not appear in the $\dot e_\gamma$ and $\dot e_V$ rows
because $f_\gamma$ and $f_V$ do not depend on $\psi$ once rotation is dropped.

\subsection*{Step 5: state-space form}

With $\vb e = [\,e_c,\ e_h,\ e_\psi,\ e_\gamma,\ e_V\,]^T$ and $\delta\vb u = [\,\delta\alpha,\ \delta\sigma\,]^T$:
\begin{equation}
  \dot{\vb e} = A_e(s)\,\vb e + B_e(s)\,\delta\vb u,
\end{equation}
\begin{equation}
A_e =
\begin{bmatrix}
0 & 0 & \sdot^* & 0 & 0\\[4pt]
-V^*\kappa\sin\gamma^* & \dfrac{V^*\sin\gamma^*}{r^*} & 0 & \dfrac{V^*}{\cos\gamma^*} & 0\\[10pt]
-\kappa\dot\psi^* & a_{\psi r}+\dfrac{\dot\psi^*}{r^*} & a_{\psi\psi} & a_{\psi\gamma}+\dot\psi^*\tan\gamma^* & a_{\psi V}-\dfrac{\dot\psi^*}{V^*}\\[10pt]
-\kappa\dot\gamma^* & a_{\gamma r}+\dfrac{\dot\gamma^*}{r^*} & 0 & a_{\gamma\gamma}+\dot\gamma^*\tan\gamma^* & a_{\gamma V}-\dfrac{\dot\gamma^*}{V^*}\\[10pt]
-\kappa\dot V^* & a_{Vr}+\dfrac{\dot V^*}{r^*} & 0 & a_{V\gamma}+\dot V^*\tan\gamma^* & a_{VV}-\dfrac{\dot V^*}{V^*}
\end{bmatrix},
\qquad
B_e =
\begin{bmatrix}
0 & 0\\[4pt] 0 & 0\\[10pt] b_{\psi\alpha} & b_{\psi\sigma}\\[10pt] b_{\gamma\alpha} & b_{\gamma\sigma}\\[10pt] b_{V\alpha} & 0
\end{bmatrix}.
\label{eq:errmodel}
\end{equation}

\section{Feedforward and why the model has no disturbance term}

At $\vb e = 0$ and $\delta\vb u = 0$, every row of the nonlinear error dynamics is exactly
zero, because the reference input $\vb u^*(s)$ reproduces the reference motion. In
Rajamani's lane-keeping model, the road curvature enters as a disturbance term
($\dot\psi_{\text{des}}$) that the feedback must reject unless a feedforward steering angle
is added. Here the planner already supplies that feedforward: $\alpha^*(s)$ and
$\sigma^*(s)$ are the inputs that fly the path exactly. The commands are
\begin{equation}
  \alpha_c = \alpha^*(s) + \delta\alpha,\qquad \sigma_c = \sigma^*(s) + \delta\sigma,
\end{equation}
and the feedback only corrects deviations.

\section{Arc length as the independent variable}

To schedule on $s$ rather than time, divide by $\sdot$. To first order,
$d\vb e/ds = (1/\sdot^*)\,\dot{\vb e}$, because the correction to $\sdot$ multiplies $\dot{\vb e}$,
which is already first order:
\begin{equation}
  \frac{d\vb e}{ds} = \frac{1}{\sdot^*}A_e(s)\,\vb e + \frac{1}{\sdot^*}B_e(s)\,\delta\vb u .
\end{equation}

\section{Structure and controllability}

\begin{itemize}
\item \textbf{Lateral chain:} $\delta\sigma \xrightarrow{\ b_{\psi\sigma}\ } e_\psi
  \xrightarrow{\ \sdot^*\ } e_c$. A double integrator from bank to crossrange, exactly like
  steering to lateral offset in Rajamani's model. $b_{\psi\sigma} = \aL\cos\sigma^*/(V^*\cos\gamma^*)$
  is largest near wings-level and vanishes at $\sigma^* = \pm90^\circ$.
\item \textbf{Vertical chain:} $(\delta\alpha,\delta\sigma) \to e_\gamma \to e_h$, with
  $b_{\gamma\sigma} = -\aL\sin\sigma^*/V^*$. Bank changes vertical lift only when the
  reference is already banked.
\item \textbf{Energy:} $\delta\alpha \xrightarrow{\ b_{V\alpha}\ } e_V$ through drag, and
  $e_\gamma \xrightarrow{\ a_{V\gamma}=-g\cos\gamma\ } e_V$ through gravity. This is what
  replaces the throttle loop of a powered vehicle.
\item \textbf{Locality.} The model is valid for small errors. Large crossrange corrections in
  entry guidance are made by \emph{bank reversals} (sign changes of $\sigma$), which are
  outside any linear model and are usually handled by logic, such as a heading-error
  corridor.
\end{itemize}

\section{What was neglected}

\begin{itemize}
\item Earth rotation and $J_2$ in the partials (\cref{sec:simplify}). Keep them by using exact
  Jacobians in code.
\item The effect of $e_c$ on latitude, which would add $a_{\psi\phi}\,(-\sin\psi^*/R_0)\,e_c$ to
  the heading row: about $10^{-7}$ rad/s per km of crossrange.
\item The difference between the heading reference at the vehicle and at the projected
  point (meridian convergence), of order $e_c\tan\phi/R_0$: about $0.005^\circ$ per km.
\end{itemize}

%======================================================================
\chapter{LQR design on the error-state model}\label{ch:lqr}

\section{Cost and weights}

\begin{equation}
  J = \int\left(\vb e^T Q\,\vb e + \delta\vb u^T R\,\delta\vb u\right)dt,\qquad
  \delta\vb u = -K(s)\,\vb e .
\end{equation}
Bryson's rule sets diagonal weights from the largest acceptable value of each quantity:
$Q_{ii} = 1/e_{i,\max}^2$ and $R_{jj} = 1/\delta u_{j,\max}^2$. Reasonable starting values:
\begin{center}
\begin{tabular}{@{}lll@{\qquad}lll@{}}
\toprule
Error & Largest acceptable & $Q_{ii}$ & Input & Largest correction & $R_{jj}$ \\
\midrule
$e_c$ & 1 km & $10^{-6}$ m$^{-2}$ & $\delta\alpha$ & $2^\circ$ & $1/(0.0349)^2$ \\
$e_h$ & 500 m & $4\times10^{-6}$ m$^{-2}$ & $\delta\sigma$ & $10^\circ$ & $1/(0.1745)^2$ \\
$e_\psi$ & $1^\circ$ & $1/(0.0175)^2$ & & & \\
$e_\gamma$ & $0.5^\circ$ & $1/(0.0087)^2$ & & & \\
$e_V$ & 20 m/s & $2.5\times10^{-3}$ (m/s)$^{-2}$ & & & \\
\bottomrule
\end{tabular}
\end{center}

\section{Computing the gains}

\paragraph{Frozen-point (gain scheduling).} At each reference node $s_k$, solve the algebraic
Riccati equation
\begin{equation}
  A_e^T P + P A_e - P B_e R^{-1} B_e^T P + Q = 0,\qquad K_k = R^{-1}B_e^T P,
\end{equation}
and interpolate $K(s)$ between nodes (\code{scipy.linalg.solve\_continuous\_are}). This is the
form that carries over directly to aircraft gain scheduling. It assumes the dynamics change
slowly relative to the closed-loop bandwidth.

\paragraph{Finite horizon.} The trajectory ends at a target, so the more principled form
integrates the differential Riccati equation backward from the end:
\begin{equation}
  -\dot P = A_e^T P + P A_e - P B_e R^{-1} B_e^T P + Q,\qquad P(t_f) = Q_f,\qquad
  K(t) = R^{-1}B_e^T P(t).
\end{equation}
In arc length, use $A_e/\sdot^*$, $B_e/\sdot^*$ and weights $Q/\sdot^*$, $R/\sdot^*$
(since $dt = ds/\sdot$), integrating backward in $s$.

\section{The control law}

\begin{align}
  \alpha_c &= \alpha^*(s) - K_{\alpha c}e_c - K_{\alpha h}e_h - K_{\alpha\psi}e_\psi - K_{\alpha\gamma}e_\gamma - K_{\alpha V}e_V, \\
  \sigma_c &= \sigma^*(s) - K_{\sigma c}e_c - K_{\sigma h}e_h - K_{\sigma\psi}e_\psi - K_{\sigma\gamma}e_\gamma - K_{\sigma V}e_V,
\end{align}
then saturate: $\alpha_{\text{lo}}(M)\le\alpha_c\le\alpha_{\text{hi}}(M)$, $|\sigma_c|\le\sigma_{\max}$,
and rate-limit if the inner loop needs it.

\section{Getting $A$ and $B$ numerically}\label{sec:casadi}

The hand-derived partials are for understanding and for sanity checks. In code, take exact
Jacobians of the full model (rotation, $J_2$, spline aerodynamics) with CasADi. The planner's
model treats $\alpha$ and $\sigma$ as states; for guidance they are inputs, so split the state:
\begin{lstlisting}
import casadi as ca
x = ca.MX.sym("x", 6)                  # [r, lon, lat, V, gamma, psi]
u = ca.MX.sym("u", 2)                  # [alpha, sigma]: inputs to the guidance loop
xdot = model.f(ca.vertcat(x, u), ca.DM.zeros(2))[:6]
AB = ca.Function("AB", [x, u], [ca.jacobian(xdot, x), ca.jacobian(xdot, u)])
A, B = AB(res.x[:6, k], res.x[6:, k])  # at node k of the plan
\end{lstlisting}
The entries of $A$ and $B$ correspond to the $a_{ij}$ and $b_{ij}$ of \cref{ch:lin}. $A_e$ and
$B_e$ then follow from \cref{eq:errmodel} using the reference rates $\dot V^*$, $\dot\gamma^*$,
$\dot\psi^*$ and the curvature $\kappa$. Alternatively, write the nonlinear error dynamics
\cref{eq:ne_c}--\cref{eq:ne_V} as a CasADi function of $(\vb e,\delta\vb u)$ and differentiate it
at zero.

\section{Checking the design}

Classical verification applies whatever the design method: break the loop at each input
($\delta\alpha$, $\delta\sigma$) of the closed loop, compute the loop transfer function, and report gain and
phase margins at frozen points along the path. Continuous-time LQR with full-state feedback
guarantees at least $60^\circ$ of phase margin and a gain margin from $-6$ dB to $+\infty$ at
the plant input, but those guarantees disappear once an estimator, actuator lag or time delay
is in the loop (Doyle, 1978), so they must be checked. Then run Monte Carlo dispersions
(density, aerodynamic coefficients, mass, initial state) in the simulator.

%======================================================================
\chapter{Inner-loop attitude control: classical and NDI}\label{ch:inner}

In a 6-DOF simulation, guidance commands $\alpha_c$, $\sigma_c$ and $\beta_c = 0$, and an
attitude control loop makes the airframe follow them using control surfaces (and reaction
control jets at low dynamic pressure). This chapter gives the models for that loop and the
structure of the two common designs. It uses \emph{forces} $L$, $D$, $Y$ (newtons) rather than
the accelerations of earlier chapters, with $\aL = L/m$.

\section{Notation for this chapter}

\begin{center}
\begin{tabular}{@{}ll@{}}
\toprule
Symbol & Meaning \\
\midrule
$p, q, r$ & body-axis roll, pitch, yaw rates \\
$\alpha, \beta$ & angle of attack, sideslip \\
$\mu$ & wind-axis bank angle (the $\sigma$ of earlier chapters) \\
$\chi$ & flight-path heading (the $\psi$ of earlier chapters) \\
$L, D, Y$ & lift, drag, side force in wind axes [N] \\
$\mathcal L, \mathcal M, \mathcal N$ & rolling, pitching, yawing moments about the center of mass [N\,m] \\
$\vb\delta$ & vector of control-surface deflections \\
\bottomrule
\end{tabular}
\end{center}
For the fast attitude dynamics, the Earth is treated as flat and non-rotating. The neglected
terms are of order $V/r$ and $\Omega$ (about $10^{-3}$ rad/s), far below attitude bandwidths
(about 1 rad/s).

\section{Rotational dynamics}

Euler's equation $I\dot{\vb\omega} + \vb\omega\times I\vb\omega = \vb M$, with the symmetric
inertia tensor
$I = \begin{bmatrix} I_{xx} & 0 & -I_{xz}\\ 0 & I_{yy} & 0\\ -I_{xz} & 0 & I_{zz}\end{bmatrix}$,
written out by component:
\begin{align}
  I_{xx}\dot p - I_{xz}\dot r &= \mathcal L + (I_{yy}-I_{zz})\,q r + I_{xz}\,p q, \\
  I_{yy}\dot q &= \mathcal M + (I_{zz}-I_{xx})\,p r + I_{xz}\,(r^2 - p^2), \\
  I_{zz}\dot r - I_{xz}\dot p &= \mathcal N + (I_{xx}-I_{yy})\,p q - I_{xz}\,q r .
\end{align}
The HL-20 data set $I_{xz}=0$, so each axis can be solved on its own:
\begin{align}
  \dot p &= \frac{1}{I_{xx}}\left[\mathcal L + (I_{yy}-I_{zz})\,q r\right], \label{eq:pdot}\\
  \dot q &= \frac{1}{I_{yy}}\left[\mathcal M + (I_{zz}-I_{xx})\,p r\right], \label{eq:qdot}\\
  \dot r &= \frac{1}{I_{zz}}\left[\mathcal N + (I_{xx}-I_{yy})\,p q\right]. \label{eq:rdot_body}
\end{align}

\section{Aerodynamic moments}

\begin{align}
  \mathcal L &= \bar q S b\left[C_{l_\beta}\beta + C_{l_p}\frac{p b}{2V} + C_{l_r}\frac{r b}{2V} + C_{l,\vb\delta}\,\vb\delta\right],\\
  \mathcal M &= \bar q S \bar c\left[C_{m_0}(\alpha,M) + C_{m_q}\frac{q\bar c}{2V} + C_{m,\vb\delta}\,\vb\delta\right],\\
  \mathcal N &= \bar q S b\left[C_{n_\beta}\beta + C_{n_p}\frac{p b}{2V} + C_{n_r}\frac{r b}{2V} + C_{n,\vb\delta}\,\vb\delta\right].
\end{align}
The \code{hgv3} aero file contains $C_m$, the sideslip derivatives $C_{Y_\beta}$, $C_{l_\beta}$,
$C_{n_\beta}$ (per \emph{degree} of $\beta$) and the damping derivatives $C_{m_q}$, $C_{l_p}$,
$C_{l_r}$, $C_{n_p}$, $C_{n_r}$. The control-surface terms $C_{\cdot,\vb\delta}$ are not in that file;
take them from \code{hgv6}'s effector model.

Throughout, split each moment into a part independent of the controls and a part linear in
them:
\begin{equation}
  \mathcal L = \mathcal L_0(\vb x) + \bar q S b\, C_{l,\vb\delta}\,\vb\delta,\quad
  \mathcal M = \mathcal M_0(\vb x) + \bar q S \bar c\, C_{m,\vb\delta}\,\vb\delta,\quad
  \mathcal N = \mathcal N_0(\vb x) + \bar q S b\, C_{n,\vb\delta}\,\vb\delta .
  \label{eq:moment_split}
\end{equation}

\section{Wind-axis attitude kinematics}

These relate the body rates to the rates of the controlled angles. They are the key
equations for both designs:
\begin{align}
  \dot\alpha &= q - \tan\beta\left(p\cos\alpha + r\sin\alpha\right)
     - \frac{L - m g\cos\gamma\cos\mu}{m V\cos\beta}, \label{eq:alphadot}\\[4pt]
  \dot\beta &= p\sin\alpha - r\cos\alpha + \frac{Y + m g\cos\gamma\sin\mu}{m V}, \label{eq:betadot}\\[4pt]
  \dot\mu &= \frac{p\cos\alpha + r\sin\alpha}{\cos\beta}
     + \frac{L\left(\tan\gamma\sin\mu + \tan\beta\right) + Y\tan\gamma\cos\mu
             - m g\cos\gamma\cos\mu\tan\beta}{m V}. \label{eq:mudot}
\end{align}
The flight-path equations in the same frame are
\begin{equation}
  \dot\gamma = \frac{L\cos\mu - Y\sin\mu - m g\cos\gamma}{m V},\qquad
  \dot\chi = \frac{L\sin\mu + Y\cos\mu}{m V\cos\gamma} .
\end{equation}
With $Y=0$ and the curved-Earth terms restored, these are the 3-DOF equations
\cref{eq:s_g,eq:s_p} with $\mu = \sigma$ and $\chi = \psi$.

\paragraph{Roll about the velocity vector.} In \cref{eq:mudot}, the body rates enter as
$p\cos\alpha + r\sin\alpha$, the roll rate about the velocity vector. At the high angles of
attack used in entry ($20^\circ$ to $40^\circ$), rolling about the body $x$ axis alone
($r=0$) would, from \cref{eq:betadot}, convert directly into sideslip at the rate
$p\sin\alpha$. Banking must therefore be done about the velocity vector, with
$r = p\tan\alpha$. Both designs below have to respect this.

\section{Classical gain-scheduled design}

\subsection{Linearization about a wings-level trim}

At each design point $(M, \bar q, \alpha^*)$ along the trajectory, linearize about
$\beta = 0$, $\mu = 0$, $p = r = 0$. The trim values of $V$ and $\gamma$ are held fixed
(they change slowly compared with attitude).

\paragraph{Longitudinal (short period).} From \cref{eq:alphadot,eq:qdot}:
\begin{align}
  \delta\dot\alpha &= \delta q - \frac{\delta L}{mV}
    = Z_\alpha\,\delta\alpha + \delta q + Z_{\delta}\,\delta\delta_e, \\
  \delta\dot q &= M_\alpha\,\delta\alpha + M_q\,\delta q + M_{\delta}\,\delta\delta_e,
\end{align}
with the dimensional derivatives
\begin{align}
  Z_\alpha &= -\frac{\bar q S C_{L_\alpha}}{m V}, &
  Z_\delta &= -\frac{\bar q S C_{L,\delta_e}}{m V}, \notag\\
  M_\alpha &= \frac{\bar q S\bar c\, C_{m_\alpha}}{I_{yy}}, &
  M_q &= \frac{\bar q S\bar c^2 C_{m_q}}{2 V I_{yy}}, &
  M_\delta &= \frac{\bar q S\bar c\, C_{m,\delta_e}}{I_{yy}} .
\end{align}
Here $\delta_e$ is whichever symmetric surface produces pitching moment. The characteristic
equation is
\begin{equation}
  s^2 - \left(Z_\alpha + M_q\right)s + \left(Z_\alpha M_q - M_\alpha\right) = 0,
  \qquad
  \omega_{sp}^2 = Z_\alpha M_q - M_\alpha,\quad
  2\zeta_{sp}\omega_{sp} = -\left(Z_\alpha + M_q\right).
\end{equation}

\paragraph{Lateral-directional.} From \cref{eq:betadot,eq:mudot,eq:pdot,eq:rdot_body}:
\begin{align}
  \delta\dot\beta &= Y_\beta\,\delta\beta + \sin\alpha^*\,\delta p - \cos\alpha^*\,\delta r
     + \frac{g\cos\gamma^*}{V}\,\delta\mu + Y_{\vb\delta}\,\delta\vb\delta, \\
  \delta\dot p &= L_\beta\,\delta\beta + L_p\,\delta p + L_r\,\delta r + L_{\vb\delta}\,\delta\vb\delta, \\
  \delta\dot r &= N_\beta\,\delta\beta + N_p\,\delta p + N_r\,\delta r + N_{\vb\delta}\,\delta\vb\delta, \\
  \delta\dot\mu &= \left(\frac{L^* - m g\cos\gamma^*}{m V} + Y_\beta\tan\gamma^*\right)\delta\beta
     + \cos\alpha^*\,\delta p + \sin\alpha^*\,\delta r + \frac{L^*\tan\gamma^*}{m V}\,\delta\mu
     + \tan\gamma^*\,Y_{\vb\delta}\,\delta\vb\delta ,
\end{align}
with
\begin{align}
  Y_\beta &= \frac{\bar q S C_{Y_\beta}}{m V}, &
  L_\beta &= \frac{\bar q S b\, C_{l_\beta}}{I_{xx}}, &
  L_p &= \frac{\bar q S b^2 C_{l_p}}{2 V I_{xx}}, &
  L_r &= \frac{\bar q S b^2 C_{l_r}}{2 V I_{xx}}, \notag\\
  & &
  N_\beta &= \frac{\bar q S b\, C_{n_\beta}}{I_{zz}}, &
  N_p &= \frac{\bar q S b^2 C_{n_p}}{2 V I_{zz}}, &
  N_r &= \frac{\bar q S b^2 C_{n_r}}{2 V I_{zz}}.
\end{align}
Convert the \code{hgv3} sideslip derivatives from per degree to per radian (multiply by
$180/\pi$) before using them here.

\subsection{Control structure}

\begin{itemize}
\item \textbf{Pitch:} an inner pitch-rate loop for damping, and an outer proportional-integral
  (PI) loop on $\alpha$ that commands it:
  $q_c = K_\alpha(\alpha_c - \alpha) + K_{I\alpha}\int(\alpha_c - \alpha)\,dt$, then
  $\delta_e = K_q(q_c - q)$.
\item \textbf{Roll:} a PI loop on $\mu$ commanding the stability-axis roll rate
  $p_{s,c} = p_c\cos\alpha + r_c\sin\alpha$, implemented as $p_c$ and $r_c = p_c\tan\alpha$ so
  the roll is about the velocity vector.
\item \textbf{Yaw:} drive $\beta$ to zero (turn coordination), with yaw-rate damping through a
  washout filter so steady turning is not opposed.
\item \textbf{Scheduling:} control power scales with $\bar q$, so gains are scheduled on
  $\bar q$ and Mach, at minimum $K \propto 1/\bar q$ to hold loop gain constant.
\item \textbf{Requirements:} at every design point, with the loop broken at each actuator,
  at least 6 dB of gain margin and $45^\circ$ of phase margin. Keep bandwidths separated:
  guidance about 0.05--0.2 rad/s, attitude about 1--3 rad/s, actuators 20 rad/s or more.
\end{itemize}

\section{Nonlinear dynamic inversion (NDI)}

\subsection{The idea, written out}

Suppose a set of outputs $\vb y$ obeys
\begin{equation}
  \dot{\vb y} = \vb f(\vb x) + G(\vb x)\,\vb v,
  \label{eq:affine}
\end{equation}
where $\vb v$ is the input, $\vb f$ collects everything the input does not affect, and $G$ is
square and invertible. Choose the input
\begin{equation}
  \vb v = G(\vb x)^{-1}\left[\vb\nu - \vb f(\vb x)\right].
\end{equation}
Substituting back into \cref{eq:affine}:
\begin{equation}
  \dot{\vb y} = \vb f(\vb x) + G(\vb x)\,G(\vb x)^{-1}\left[\vb\nu - \vb f(\vb x)\right] = \vb\nu .
\end{equation}
The nonlinear system now behaves as a set of decoupled integrators driven by the
\emph{pseudo-control} $\vb\nu$. Choose $\vb\nu$ as a simple linear law, for example
\begin{equation}
  \nu_i = \dot y_{i,c} + K_{P,i}\left(y_{i,c} - y_i\right) + K_{I,i}\int\left(y_{i,c} - y_i\right)dt,
\end{equation}
which gives each channel first-order (or, with the integral, second-order) linear error
dynamics set by the gains. Conditions: $G$ must be invertible along the trajectory, and the
\emph{internal dynamics}, the states not directly controlled ($V$, $\gamma$, $\chi$ here), must
be stable or handled by another loop (guidance).

For the HL-20, the controlled angles $(\alpha,\beta,\mu)$ depend on the body rates, and the
body rates depend on the control surfaces. So the inversion is done in two stages, using
time-scale separation: the rates respond much faster than the angles.

\subsection{Slow loop: angles to body rates}

From \cref{eq:alphadot,eq:betadot,eq:mudot}, the angle rates are affine in $(p, q, r)$:
\begin{equation}
  \begin{bmatrix}\dot\alpha\\ \dot\beta\\ \dot\mu\end{bmatrix}
  = \begin{bmatrix} f_\alpha\\ f_\beta\\ f_\mu \end{bmatrix}
  + G_s(\alpha,\beta)\begin{bmatrix} p\\ q\\ r\end{bmatrix},
\end{equation}
\begin{align}
  f_\alpha &= -\frac{L - m g\cos\gamma\cos\mu}{m V\cos\beta}, \\
  f_\beta &= \frac{Y + m g\cos\gamma\sin\mu}{m V}, \\
  f_\mu &= \frac{L\left(\tan\gamma\sin\mu + \tan\beta\right) + Y\tan\gamma\cos\mu
            - m g\cos\gamma\cos\mu\tan\beta}{m V}, \\[4pt]
  G_s &= \begin{bmatrix}
    -\tan\beta\cos\alpha & 1 & -\tan\beta\sin\alpha\\
    \sin\alpha & 0 & -\cos\alpha\\
    \dfrac{\cos\alpha}{\cos\beta} & 0 & \dfrac{\sin\alpha}{\cos\beta}
  \end{bmatrix},\qquad
  \det G_s = -\frac{1}{\cos\beta}.
\end{align}
$G_s$ is invertible for all $|\beta|<90^\circ$, with
\begin{equation}
  G_s^{-1} = \begin{bmatrix}
    0 & \sin\alpha & \cos\alpha\cos\beta\\
    1 & 0 & \sin\beta\\
    0 & -\cos\alpha & \sin\alpha\cos\beta
  \end{bmatrix}.
\end{equation}
With desired angle rates $\nu_\alpha$, $\nu_\beta$, $\nu_\mu$, the commanded body rates are,
written out:
\begin{align}
  p_c &= \sin\alpha\left(\nu_\beta - f_\beta\right) + \cos\alpha\cos\beta\left(\nu_\mu - f_\mu\right), \\
  q_c &= \left(\nu_\alpha - f_\alpha\right) + \sin\beta\left(\nu_\mu - f_\mu\right), \\
  r_c &= -\cos\alpha\left(\nu_\beta - f_\beta\right) + \sin\alpha\cos\beta\left(\nu_\mu - f_\mu\right).
\end{align}
With $\beta \approx 0$, a pure bank command gives $p_c = \cos\alpha\,\nu_\mu'$ and
$r_c = \sin\alpha\,\nu_\mu'$, so $r_c/p_c = \tan\alpha$: the inversion automatically rolls
about the velocity vector. The desired rates are simple linear laws:
\begin{equation}
  \nu_\alpha = \omega_\alpha(\alpha_c - \alpha),\qquad
  \nu_\beta = \omega_\beta(0 - \beta),\qquad
  \nu_\mu = \omega_\mu(\mu_c - \mu),
\end{equation}
optionally with integral terms. $L$ and $Y$ in $f_\alpha$, $f_\beta$, $f_\mu$ come from the
aerodynamic model at the current state.

\subsection{Fast loop: body rates to control moments}

From \cref{eq:pdot,eq:qdot,eq:rdot_body} with the split \cref{eq:moment_split}:
\begin{align}
  \dot p &= \frac{(I_{yy}-I_{zz})\,q r + \mathcal L_0}{I_{xx}} + \frac{\bar q S b}{I_{xx}}C_{l,\vb\delta}\,\vb\delta, \\
  \dot q &= \frac{(I_{zz}-I_{xx})\,p r + \mathcal M_0}{I_{yy}} + \frac{\bar q S \bar c}{I_{yy}}C_{m,\vb\delta}\,\vb\delta, \\
  \dot r &= \frac{(I_{xx}-I_{yy})\,p q + \mathcal N_0}{I_{zz}} + \frac{\bar q S b}{I_{zz}}C_{n,\vb\delta}\,\vb\delta .
\end{align}
With desired rate derivatives $\nu_p = \omega_p(p_c - p)$, $\nu_q = \omega_q(q_c - q)$,
$\nu_r = \omega_r(r_c - r)$, solve each axis for the moment the surfaces must produce:
\begin{align}
  \mathcal L_c &= I_{xx}\,\nu_p - (I_{yy}-I_{zz})\,q r - \mathcal L_0, \\
  \mathcal M_c &= I_{yy}\,\nu_q - (I_{zz}-I_{xx})\,p r - \mathcal M_0, \\
  \mathcal N_c &= I_{zz}\,\nu_r - (I_{xx}-I_{yy})\,p q - \mathcal N_0 .
\end{align}

\paragraph{Control allocation.} The surfaces must satisfy
\begin{equation}
  \bar q S\begin{bmatrix} b\,C_{l,\vb\delta}\\ \bar c\,C_{m,\vb\delta}\\ b\,C_{n,\vb\delta}\end{bmatrix}\vb\delta
  \equiv B_\delta\,\vb\delta
  = \begin{bmatrix}\mathcal L_c\\ \mathcal M_c\\ \mathcal N_c\end{bmatrix}.
\end{equation}
With exactly three independent effectors, $\vb\delta = B_\delta^{-1}[\mathcal L_c,\mathcal M_c,\mathcal N_c]^T$.
With more than three, use a weighted pseudo-inverse,
$\vb\delta = W^{-1}B_\delta^T\left(B_\delta W^{-1}B_\delta^T\right)^{-1}[\mathcal L_c,\mathcal M_c,\mathcal N_c]^T$,
then clip to deflection limits (or use a constrained allocator).

\subsection{Incremental NDI (INDI)}

NDI needs accurate $\mathcal L_0$, $\mathcal M_0$, $\mathcal N_0$. INDI avoids most of that
model. Write the rate dynamics as $\dot{\vb\omega} = \vb f_\omega(\vb x) + G_\omega\vb\delta$
with $G_\omega = \mathrm{diag}(I_{xx},I_{yy},I_{zz})^{-1}B_\delta$. At the previous sample,
the measured angular acceleration is $\dot{\vb\omega}_0 = \vb f_\omega(\vb x_0) + G_\omega\vb\delta_0$.
Subtracting, and assuming $\vb f_\omega$ changes little over one sample:
\begin{equation}
  \dot{\vb\omega} - \dot{\vb\omega}_0 \approx G_\omega\left(\vb\delta - \vb\delta_0\right)
  \quad\Rightarrow\quad
  \vb\delta = \vb\delta_0 + G_\omega^{-1}\left(\vb\nu_\omega - \dot{\vb\omega}_0\right).
\end{equation}
Only the control effectiveness $G_\omega$ is needed. The costs are an angular acceleration
estimate (filtered differentiation of the rate gyros), knowledge of the actual actuator
positions $\vb\delta_0$, and matched filtering of the two.

\subsection{Practical points}

\begin{itemize}
\item \textbf{Bandwidth separation:} $\omega_p,\omega_q,\omega_r$ about 3--5 times
  $\omega_\alpha,\omega_\beta,\omega_\mu$, which are about 3--5 times the guidance
  bandwidth. Actuators faster still.
\item \textbf{Saturation:} when surfaces saturate the inversion is no longer exact; use
  anti-windup on integrators and consider pseudo-control hedging.
\item \textbf{Verification:} NDI's nominal loop is an integrator ($\nu/s$), which looks
  perfectly robust. Real margins come from actuator dynamics, delays and model error, so
  linearize the closed loop with those included, break it at the actuators, and check
  margins exactly as for the classical design. Repeat with perturbed control effectiveness and
  moment coefficients.
\end{itemize}

\section{Comparison}

\begin{center}
\renewcommand{\arraystretch}{1.3}
\begin{tabular}{@{}>{\raggedright\arraybackslash}p{0.2\textwidth}>{\raggedright\arraybackslash}p{0.24\textwidth}>{\raggedright\arraybackslash}p{0.24\textwidth}>{\raggedright\arraybackslash}p{0.24\textwidth}@{}}
\toprule
& Classical, gain-scheduled & NDI & INDI \\
\midrule
Model needed & linear models at design points & full moment model and effectiveness & effectiveness only \\
Scheduling & explicit gain tables & built in & built in \\
Sensitivity to aero error & moderate (margins) & high & low \\
Extra sensors & none & none & angular acceleration, actuator positions \\
Certification history & long (Shuttle, most aircraft) & established (X-38, modern fighters) & growing (research, eVTOL) \\
\bottomrule
\end{tabular}
\end{center}

%======================================================================
\appendix
\chapter{Verification record}\label{app:verify}

\begin{center}
\renewcommand{\arraystretch}{1.25}
\begin{tabular}{@{}>{\raggedright\arraybackslash}p{0.36\textwidth}>{\raggedright\arraybackslash}p{0.6\textwidth}@{}}
\toprule
Item & How it was checked \\
\midrule
Local basis derivatives, $\vb\omega_{LE}$, kinematics, transport terms (\cref{ch:eom}) &
  Symbolic (sympy): every intermediate expression reduces to zero difference. \\
Force projections and assembled equations \cref{eq:full_r}--\cref{eq:full_psi} &
  Symbolic: derived from \cref{eq:rotating} and compared with the expressions in
  \code{hl20\_model.py}: identical. \\
$J_2$ gravity \cref{eq:gravity} & Numerical against \code{hgv3.gravity\_vector} at 50 random
  points: relative error below $10^{-15}$. \\
Model against simulator & Open-loop $(\alpha,\sigma)$ from Mach 15 for 800 s at three bank
  settings: altitude within 10 m, speed within 0.1 m/s, along-track drift 0.8 m/s. \\
Partial derivatives (\cref{ch:lin}) & All 34 hand expressions against symbolic
  differentiation with generic $C_L$, $C_D$, $\rho$, $a_s$: relative error below
  $10^{-14}$; all entries listed as zero confirmed zero. \\
Simplification error (\cref{sec:simplify}) & CasADi Jacobians of the full and simplified
  model along a planned trajectory. \\
Error-state model \cref{eq:errmodel} & Symbolic Jacobian of the nonlinear error dynamics with
  generic $f_V$, $f_\gamma$, $f_\psi$: matches entry by entry; error rates vanish at zero
  error. \\
CasADi snippet (\cref{sec:casadi}) & Run against \code{hl20\_model.py} on a planned
  trajectory. \\
Euler equations with $I_{xz}$ & Symbolic: identical to $I\dot{\vb\omega}+\vb\omega\times I\vb\omega=\vb M$. \\
Wind-axis kinematics \cref{eq:alphadot}--\cref{eq:mudot} and $\dot\gamma$, $\dot\chi$ &
  Numerical against a rigid-body model (finite difference of angles extracted from the
  attitude matrix and velocity) at 200 random states with sideslip and side force: agreement
  to $10^{-10}$. \\
$G_s^{-1}$ and the wings-level linearizations & Symbolic. \\
\bottomrule
\end{tabular}
\end{center}

\end{document}
